% year: תשפ"ד
% type: מבחן
% semester: א
% moed: ב
% number: 4
% points: 14
% categories: taylor
% source: exams/מבחנים/תשפד/מועד ב תשפד.pdf

\begin{question}
העריכו בעזרת קירוב טיילור מסדר שני את $\sqrt[3]{e}$, ותנו חסם לשגיאה.
\end{question}

\begin{hint}
השתמשו בפולינום מקלורן של $e^x$ מסדר $2$ בנקודה $x=\frac13$, ובשארית בצורת לגרנז'.
\end{hint}

\begin{solution}
נשתמש בפונקציה $f(x)=e^x$ סביב $x_0=0$, ונציב $x=\frac13$, שכן $\sqrt[3]{e}=e^{1/3}$.

\textbf{פולינום טיילור.} לכל $n$, $f^{(n)}(x)=e^x$ ולכן $f^{(n)}(0)=1$. פולינום מקלורן מסדר $2$:
\[ P_2(x)=1+x+\frac{x^2}{2}. \]
לכן
\[ \sqrt[3]{e}\approx P_2\!\left(\tfrac13\right)=1+\frac13+\frac1{18}=\frac{25}{18}\approx1.3889. \]

\textbf{חסם לשגיאה.} לפי משפט טיילור עם שארית לגרנז', קיימת $c$ בין $0$ ל-$\frac13$ כך ש-
\[ R_2\!\left(\tfrac13\right)=e^{1/3}-P_2\!\left(\tfrac13\right)=\frac{f'''(c)}{3!}\left(\frac13\right)^3=\frac{e^c}{6\cdot27}=\frac{e^c}{162}. \]
מאחר ש-$e^x$ עולה ו-$0<c<\frac13$, מתקיים $e^c<e^{1/3}$. כיוון ש-$e<3<\frac{27}{8}=1.5^3$, נקבל $e^{1/3}<1.5$. לכן
\[ 0<R_2\!\left(\tfrac13\right)<\frac{1.5}{162}=\frac1{108}\approx0.0093. \]
(אפשר גם להשתמש בחסם הגס יותר $e^c<e<3$ ולקבל שגיאה קטנה מ-$\frac{3}{162}=\frac1{54}$.)

\textbf{תשובה:} $\sqrt[3]{e}\approx\frac{25}{18}\approx1.3889$, והשגיאה חיובית וקטנה מ-$\frac1{108}$. (לבדיקה: $\sqrt[3]{e}=1.39561\ldots$, והשגיאה בפועל היא כ-$0.0067$.)
\end{solution}
